\documentclass[../NDCNNAI_main.tex]{subfiles}
\graphicspath{{\subfix{../fig/}}}

\begin{document}

\subsection{Четыре классические динамические модели (Нарендра и Партасарати) для идентификации и управления}

В своей основополагающей работе 1990 года К.С. Нарандра и К. Партасарати предложили четыре базовые нелинейные модели для идентификации и управления дискретными системами <<вход-состояние-выход>>. Для SISO-систем они записываются следующим образом:
\begin{enumerate}[label=М \Roman*:]
    \item $y_p(k+1) = \sum_{i=0}^{n-1} \alpha_i y_p(k-i) + g\big(u(k), u(k-1), \dots, u(k-m+1)\big)$,
    \item $y_p(k+1) = f\big(y_p(k), y_p(k-1), \dots, y_p(k-n+1)\big) + \sum_{j=0}^{m-1} \beta_j u(k-j)$,
    \item $y_p(k+1) = f\big(y_p(k), y_p(k-1), \dots, y_p(k-n+1)\big) + g\big(u(k), u(k-1) \dots, u(k-m+1)\big)$,
    \item $y_p(k+1) = f\big(y_p(k), \dots, y_p(k-n+1), u(k), \dots, u(k-m+1)\big)$.
\end{enumerate}

Неизвестные нелинейности \(f\) и \(g\) аппроксимируются многослойными перцептронами (MLP). Модель I содержит нелинейность только по входу, модель II --- только по выходу, модель III разделяет нелинейности, а модель IV является общей НАРСС (NARX) моделью.

Эти модели мотивированы аналогичными структурами для линейных систем (ARX) и позволяют естественным образом встроить нейронные сети в динамический контекст. В оригинальной статье приведены примеры идентификации для каждой модели (например, система \(y_p(k+1) = 0.3y_p(k) + 0.6y_p(k-1) + f(u(k))\) для модели I \cite{Narendra1990}).

\subsection{Представление моделей как частных случаев GTN}

Обобщённые трансляционные сети (GTN) представляют собой композицию линейного оператора сэмплирования \(L\) и нелинейного оператора \(N\): \(P \approx N \circ L\). Модели Нарандры–Партасарати естественно вписываются в эту парадигму.

\begin{itemize}
    \item \textbf{Модели I, II, IV (NARX-типа)} непосредственно соответствуют схеме \(N \circ L\). Линейный оператор \(L_D\) формирует вектор-регрессор из запомненных значений входа и выхода:
    \[
    \phi(k) = \big( y(k-1), \dots, y(k-n_y),\; u(k-1), \dots, u(k-n_u) \big) \in \mathbb{R}^{D}.
    \]
    Нейронная сеть \(\Phi\) выступает в роли нелинейного оператора \(N\), отображая \(\phi(k)\) в предсказание \(\hat{y}(k|k-1)\). Таким образом, для модели IV (и её частных случаев I, II) имеем \(\hat{y}(k+1) = \Phi(\phi(k)) = (N \circ L_D)(u, y, k)\).

    \item \textbf{Модель III} может быть представлена как комбинация двух MLP: \(N_f\) для аппроксимации \(f\) и \(N_g\) для аппроксимации \(g\), работающих параллельно на разных частях регрессора.
    
    \item \textbf{Модель II (State-Space)} также допускает операторную интерпретацию, но в более общем виде. Система
    \[
    x(k+1) = \Psi(x(k), u(k)), \quad \hat{y}(k) = \Theta(x(k)),
    \]
    где \(\Psi\) и \(\Theta\) --- MLP, представляет собой рекуррентную нейронную сеть. Оператор сдвига во времени \(S\) действует как обобщённая трансляция на пространстве состояний \(x\). Линейное отображение, восстанавливающее состояние из прошлых данных (например, через наблюдатель или энкодер \(E(\phi(k_0))\)), можно рассматривать как часть оператора \(L\), а \(\Theta\) --- как оператор \(N\).
\end{itemize}

Все четыре модели являются частными случаями или вариациями общей архитектуры $L$--$N$, что обеспечивает единый подход к их анализу на основе теорем об универсальной аппроксимации операторов.

\subsection{Идея замыкания контура обратной связи, замечания про устойчивость и штрафы на основе якобиана/ISS}

При использовании обученной модели в контуре управления (например, в схеме косвенного адаптивного управления или MPC) модель переводится в \emph{параллельный (free-run)} режим, где её собственные предсказания \(\hat{y}\) используются для формирования регрессора на следующих шагах. Это замыкает контур обратной связи и может привести к неустойчивости, даже если исходная система устойчива.

Критически важным становится обеспечение \emph{устойчивости} модели. Для моделей в пространстве состояний (модель II) достаточным условием глобальной устойчивости является \emph{сжатие} (contraction):
\[
\sup_{x, u} \| \partial_x \Psi(x, u, \theta) \| \leq \lambda < 1.
\]
Более общим условием является \emph{устойчивость вход-состояние (ISS)}: существует функция Ляпунова \(V(x)\) и функции сравнения \(\underline{\alpha}, \overline{\alpha}, \alpha, \sigma\) такие, что
\[
\underline{\alpha}(\|x\|) \leq V(x) \leq \overline{\alpha}(\|x\|), \quad V(\Psi(x, u)) - V(x) \leq -\alpha(\|x\|) + \sigma(\|u\|),
\]
что гарантирует ограниченность состояния при ограниченном входе.

Для принудительного выполнения этих условий в процессе обучения в функционал потерь добавляют регуляризационные штрафы на якобианы:
\[
\mathcal{R}(\theta) = \lambda_x \sum_k \|\partial_x \Psi(x(k), u(k))\|_F^2 + \lambda_y \sum_k \|\partial_x \Theta(x(k))\|_F^2 + \lambda_{sn} \sum_\ell (\sigma_{\max}(W_\ell) - \tau)_+^2.
\]

Такая регуляризация не только улучшает устойчивость замкнутого контура, но и стабилизирует вычисление градиентов в BPTT, смягчая проблемы исчезающих и взрывающихся градиентов, возникающих при перемножении матриц Якоби \(\prod_j A(j)\).

Для моделей NARX-типа (I, III, IV) в параллельном режиме также возникают длинные цепочки вычислений, требующие усечённого BPTT с горизонтом \(H\) (обычно \(H \in [5, 20]\)).

\subsection{Обучающие протоколы и связь с операторным подходом Чен--Чена}

Обучение динамических моделей может проводиться в двух основных режимах:

\begin{enumerate}
    \item \textbf{Последовательно-параллельный (Series-Parallel)} или \emph{teacher forcing}: в процессе обучения в регрессор подставляются измеренные значения выхода \(y(k-i)\). Это разрывает рекуррентные связи и превращает задачу в обучение статического отображения, что позволяет использовать стандартное обратное распространение. Соответствующая функция потерь --- \emph{одношаговая}:
    \[
    J_{\text{sp}}(\theta) = \sum_k \| y(k) - \hat{y}(k|k-1; \theta) \|^2.
    \]
    Этот режим обеспечивает устойчивое обучение (нет накопления ошибки), но может приводить к рассогласованию при переходе в параллельный режим \cite{Narendra1990}.
    
    \item \textbf{Параллельный (Parallel)} или \emph{free-run}: модель работает автономно, используя собственные предсказания. Обучение требует BPTT по нескольким шагам (горизонт \(H\)) и минимизации \emph{многошаговой} ошибки:
    \[
    J_{\text{par},H}(\theta) = \sum_{k} \sum_{h=1}^{H} \| y(k+h) - \hat{y}_{h\text{-step}}(k; \theta) \|^2.
    \]
    Этот режим лучше соответствует реальной эксплуатации модели в контуре управления, но сложнее в обучении из-за проблем с градиентами. Часто используют гибридный подход: начинают с \(J_{\text{sp}}\), затем дообучают на \(J_{\text{par},H}\) \cite{lec04}.
\end{enumerate}

Операторная теорема Чен--Чена об \hyperref[thm:chen_chen_oper]{универсальной аппроксимации нелинейных операторов} предоставляет теоретическое обоснование подхода $L$--$N$. Она утверждает, что используя фиксированные ridge-функции \(g(w_i \cdot x + \theta_i)\) и настраивая только линейные веса \(c_i(f)\) на выходе, можно равномерно аппроксимировать любой непрерывный оператор на компактном множестве входных сигналов.

Модели Нарандры–Партасарати реализуют эту идею на практике: оператор \(L\) (формирование регрессора из задержек) выполняет линейное сэмплирование, а MLP \(\Phi\) реализует нелинейное считывание. Благодаря теореме о плотности MLP в пространстве непрерывных функций, \(\Phi\) может сколь угодно точно аппроксимировать требуемую нелинейность \(F^*\).

Таким образом, модели I–IV представляют собой исторически значимый мост между классической теорией идентификации нелинейных систем и современными операторными подходами на основе нейронных сетей. Их структурная простота, ясная интерпретация в рамках GTN и наличие эффективных алгоритмов обучения (статическое и динамическое обратное распространение) делают их фундаментом для многих современных методов нейросетевого моделирования динамических систем.

Первая часть лекции основана на работе \cite{Mhaskar1997}, а вторая на \cite{Narendra1990}.

\end{document}